\documentclass[11pt,a4paper]{article}
\usepackage[T1]{fontenc}
\usepackage[utf8]{inputenc}
\usepackage{lmodern}
\usepackage[margin=27mm]{geometry}
\usepackage{amsmath,amssymb,amsthm,mathtools}
\usepackage{microtype,booktabs,longtable,array,enumitem,xcolor}
\usepackage[colorlinks=true,linkcolor=blue!45!black,citecolor=blue!45!black,urlcolor=blue!50!black,breaklinks=true]{hyperref}
\usepackage{fancyhdr}
\pagestyle{fancy}\fancyhf{}\fancyhead[L]{\small Keio / Research manuscript}\fancyhead[R]{\small September 2026}\fancyfoot[C]{\thepage}
\setlength{\headheight}{14pt}
\setlength{\emergencystretch}{3em}
\setlength{\parskip}{0.35em}
\setlist{itemsep=2pt,topsep=4pt}
\newtheorem{theorem}{Theorem}[section]
\newtheorem{lemma}[theorem]{Lemma}
\newtheorem{proposition}[theorem]{Proposition}
\newtheorem{corollary}[theorem]{Corollary}
\newtheorem{conjecture}[theorem]{Conjecture}
\newtheorem{criterion}[theorem]{Criterion}
\theoremstyle{definition}\newtheorem{definition}[theorem]{Definition}
\newtheorem{example}[theorem]{Example}
\theoremstyle{remark}\newtheorem{remark}[theorem]{Remark}
\newtheorem{audit}[theorem]{Proof-status note}
\newcommand{\Q}{\mathbb Q}\newcommand{\Z}{\mathbb Z}\newcommand{\N}{\mathbb N}\newcommand{\R}{\mathbb R}\newcommand{\C}{\mathbb C}
\newcommand{\Qbar}{\overline{\mathbb Q}}
\DeclareMathOperator{\Tr}{Tr}\DeclareMathOperator{\Norm}{N}\DeclareMathOperator{\trdeg}{trdeg}\DeclareMathOperator{\ord}{ord}\DeclareMathOperator{\lcm}{lcm}
\DeclareMathOperator{\Ein}{Ein}\DeclareMathOperator{\Ei}{Ei}\DeclareMathOperator{\Ci}{Ci}\DeclareMathOperator{\Chi}{Chi}\DeclareMathOperator{\Li}{Li}
\newcommand{\floor}[1]{\left\lfloor #1\right\rfloor}\newcommand{\ceil}[1]{\left\lceil #1\right\rceil}
\newcommand{\sourceT}[1]{\textup{[source T#1]}}
\author{Keio}
\date{5 September 2026\\\small Standalone research draft; source audit identifiers retained}

\title{Minimal Recurrence Orders and Integral Solution Lattices\newline for Franel Sequences}
\newcommand{\CT}{\operatorname{CT}}
\newcommand{\supp}{\operatorname{supp}}
\newcommand{\vp}{v_p}
\begin{document}
\maketitle
\begin{abstract}
The generalized Franel sequence is $A_s(n)=\sum_{k=0}^{n}\binom nk^s$.
We develop two complementary rigidity arguments. The analytic argument represents
$A_s(n)$ as a moment of a real Laurent polynomial on a torus. We identify every
nonzero critical point, verify that its Hessian is definite, and prove that each
distinct nonzero critical value gives a genuine singularity of the continued
moment generating function. A recurrence of shift order $r$ allows at most $r$
such singularities. Together with Stoll's existence theorem this yields the
order $\lfloor(s+1)/2\rfloor$. The arithmetic argument concerns the classical
cubic Franel recurrence. Its discrete Abel identity forces the Casoratian of two
integral solutions to vanish, giving the complete lattice $\Z A_3$. The rational
solutions with a common denominator are exactly $\Q A_3$. We also give a stronger
local statement: any other rational solution has unbounded denominators at
every odd prime. The two arguments are independent, and the paper records their
hypotheses and provenance separately. These are research-draft proofs, not a
claim of completed external peer review or independently established priority.
\end{abstract}
\noindent\textbf{Keywords.} Franel numbers; polynomial recurrences; Cauchy transforms;
Morse critical points; Casoratian; integral solutions.

\section{Two rigidity questions}
For an integer $s\ge1$, put
\begin{equation}\label{eq:defA}
 A_s(n)=\sum_{k=0}^{n}\binom nk^s,\qquad n\ge0.
\end{equation}
The parameter $s$ specifies a sequence; the variable $n$ is its index. A linear
recurrence in this paper always means an identity in the variable $n$. We use
\emph{shift order}, rather than differential order or the number of nonzero
terms. These quantities must not be interchanged.

\begin{definition}\label{def:order}
A nonzero recurrence of shift order $r$ for a sequence $a$ is an identity
\begin{equation}\label{eq:recdef}
 \sum_{j=0}^{r}p_j(n)a_{n+j}=0
\end{equation}
for every sufficiently large integer $n$, where $p_j\in\Q(n)$ and
$p_0p_r\ne0$. Its shift order is the difference between the largest and smallest
shifts after zero endpoint coefficients have been removed. The minimum such
order is denoted by $\operatorname{ord}_{\Q(n)}a$.
\end{definition}

Stoll's theorem gives a recurrence for \eqref{eq:defA} of shift order at most
$\lfloor(s+1)/2\rfloor$ \cite{Stoll}. The distinction between this upper bound
and minimality is explicit in Straub--Zudilin \cite{SZ}. Their lower bound is
proved for telescoping recurrences based on the binomial-power summand. The
unrestricted lower bound is a separate question. The argument below concerns
an arbitrary recurrence \eqref{eq:recdef}; no telescoping certificate is assumed.

There is another question even when a specific recurrence is fixed. Write
$S_n=A_3(n)$. The classical Franel equation is
\begin{equation}\label{eq:franel}
 (n+1)^2 f_{n+1}=(7n^2+7n+2)f_n+8n^2f_{n-1},\qquad n\ge1.
\end{equation}
Its rational solution space has dimension two, because $f_0,f_1$ are free.
Its integral solution set need not have rank two. Tonien's conclusion
\cite{Tonien} asks whether the space contains integral sequences other than
multiples of $S$. We answer that question without changing the recurrence or
its starting index.

\begin{theorem}[Minimal order; \sourceT{39--40}]\label{thm:mainorder}
For every integer $s\ge1$,
\[
 \operatorname{ord}_{\Q(n)} A_s=\left\lfloor\frac{s+1}{2}\right\rfloor.
\]
In fact the same lower bound holds for recurrences with coefficients in
$\C(n)$ that are valid for all sufficiently large $n$.
\end{theorem}

\begin{theorem}[Integral lattice; \sourceT{58}]\label{thm:mainlattice}
Let $f=(f_n)_{n\ge0}$ solve \eqref{eq:franel}. Then
\[
 f_n\in\Z\text{ for all }n\ge0
 \quad\Longleftrightarrow\quad f=aS\text{ for some }a\in\Z.
\]
Among rational solutions, those for which one integer $D\ge1$ satisfies
$Df_n\in\Z$ for all $n$ are precisely the sequences $qS$, $q\in\Q$.
\end{theorem}

These are independent statements. The first does not classify rational
solutions of a recurrence, and the second does not prove minimality for all
$s$. Sections~\ref{sec:budget}--\ref{sec:orderproof} give the analytic argument.
Sections~\ref{sec:abel}--\ref{sec:padic} give the arithmetic argument, including
an explicit prime-adic strengthening. In particular, the lattice theorem can
be read and checked without accepting any assertion about analytic continuation.

\section{Recurrence width and singularities}\label{sec:budget}
We first record the elementary passage from recurrence coefficients to an
ordinary differential equation. This passage and the resulting singularity
budget are standard tools; the issue in the Franel application is to prove
that the candidate singularities are actually present.

\begin{lemma}[Singularity budget]\label{lem:budget}
Suppose $F(t)=\sum_{n\ge0}a_nt^n$ has a positive radius of convergence, and
$a$ satisfies a recurrence of shift order $r$ over $\C(n)$. There is a nonzero
polynomial $H(t)$ of degree at most $r$ such that every nonzero finite point at
which a branch continued from the germ of $F$ fails to be holomorphic is a zero
of $H$.
\end{lemma}
\begin{proof}
Clear denominators and discard finitely many initial indices. We can assume
$p_j\in\C[n]$, $p_0p_r\ne0$, and that \eqref{eq:recdef} holds for $n\ge N$.
Let $\vartheta=t\,d/dt$. Multiplying by $t^{n+r}$ and summing gives
\begin{equation}\label{eq:operator}
 LF=R(t),\qquad
 L=\sum_{j=0}^{r}t^{r-j}p_j(\vartheta-j),
\end{equation}
where $R$ is a polynomial arising only from the omitted initial indices.
Indeed, after putting $k=n+j$, the $j$th summand becomes
$t^{r-j}\sum_{k\ge N+j}p_j(k-j)a_kt^k$.

Let $D=\max_j\deg p_j$ and $c_j=[n^D]p_j(n)$. The coefficient of the highest
ordinary derivative $d^D/dt^D$ in $L$ is
\begin{equation}\label{eq:Hbudget}
 t^D H(t),\qquad H(t)=\sum_{j=0}^{r}c_jt^{r-j}.
\end{equation}
At least one $c_j$ is nonzero, and the powers of $t$ in this last sum are
distinct. Consequently $H\ne0$ and $\deg H\le r$; cancellation cannot destroy
the differential operator's highest derivative.

If $t_0\ne0$ and $H(t_0)\ne0$, equation \eqref{eq:operator} is an ordinary
linear differential equation near $t_0$. Every local solution on a punctured
subdomain extends holomorphically through $t_0$: write the equation as a
first-order analytic system and use its analytic fundamental matrix. This
applies to every branch obtained by continuation of $F$, since the differential
equation continues with the branch. When $D=0$, the conclusion follows directly
from $F=R/H$. Thus all the specified singularities lie among the roots of $H$.
\end{proof}

\begin{remark}
The degree $D$ of the recurrence coefficients is unrestricted. It is not the
quantity bounded below here. Even a recurrence with extremely large polynomial
coefficients has at most $r$ distinct nonzero finite singularities in the sense
of Lemma~\ref{lem:budget}. Initial exceptions change $R$, but not this conclusion.
Apparent zeros of $H$ are allowed; the lemma gives an upper bound on genuine
singularities, not an equality.
\end{remark}

\section{Definite critical points are visible}\label{sec:visibility}
The following statement supplies the analytic step used in the application.
It explicitly excludes saddle points on the critical fiber and permits a
critical value in the interior of the global range.

\begin{definition}
Let $M$ be a compact real-analytic manifold without boundary, of dimension
$d\ge1$, with a positive analytic density $d\mu$. Let $\Phi:M\to\R$ be analytic.
A value $v$ is called a \emph{same-side definite value} if the critical points
in $\Phi^{-1}(v)$ form a finite nonempty set, every such critical point is
nondegenerate, and all their Hessians are positive definite or all are negative
definite. Regular points in the same fiber are permitted.
\end{definition}

\begin{lemma}[Local density at a definite value]\label{lem:density}
Under these hypotheses, near a same-side definite value $v$ the pushforward
measure $\nu=\Phi_*\mu$ is absolutely continuous. Its density has the form
\begin{equation}\label{eq:densityform}
 \rho(v+t)=r(t)+\mathbf1_{\{\epsilon t>0\}}
             (\epsilon t)^{d/2-1}h(\epsilon t),
\end{equation}
where $r$ is $C^\infty$, $h$ is real analytic near zero with $h(0)>0$, and
$\epsilon=1$ for minima, $\epsilon=-1$ for maxima. In particular $\rho$ does not
agree almost everywhere near $v$ with a $C^\infty$ function.
\end{lemma}
\begin{proof}
Choose disjoint coordinate neighborhoods of the finitely many critical points.
By the analytic Morse lemma, the function in each neighborhood takes the form
\[
 \Phi(y)=v+\epsilon(y_1^2+\cdots+y_d^2),
\]
and the density becomes $g(y)\,dy$ with $g$ analytic and $g(0)>0$.
Choose a smooth nonnegative cutoff equal to one on a smaller neighborhood.
For small $w>0$, the entire sphere $|y|=\sqrt w$ lies in that smaller
neighborhood, so polar coordinates give its exact contribution to the density:
\begin{equation}\label{eq:spherical}
 \frac12 w^{d/2-1}\int_{\mathbb S^{d-1}}g(\sqrt w\,\omega)\,d\omega.
\end{equation}
The spherical average has a convergent power series in $w$. All odd-degree
Taylor monomials of $g$ integrate to zero, and the even-degree terms give
integer powers of $w$. Its constant term is positive. Thus each critical point
contributes the indicated one-sided power with a positive analytic leading
coefficient. All points use the same side $\epsilon$, so those coefficients add.

The complement of the smaller critical neighborhoods is compact. On its part
mapping to a sufficiently small interval about $v$, the gradient is nonzero.
Otherwise critical points outside those neighborhoods would accumulate on the
fiber $\Phi^{-1}(v)$, contrary to the choice of neighborhoods and compactness.
In submersion coordinates, pushforward of a smooth compactly supported density
is $C^\infty$: integrate the smooth transverse coordinates, differentiating
under the integral. A finite partition of unity therefore gives the function
$r$ in \eqref{eq:densityform}.

For even $d$, the exponent $d/2-1$ is an integer $k\ge0$. The one-sided term
has a nonzero jump in its $k$th derivative. For odd $d\ge3$, it has a
half-integral exponent and fails to be smooth at a finite derivative order.
For $d=1$ it is unbounded like $|t|^{-1/2}$, but is locally integrable.
Multiplication by $h$ with $h(0)>0$ and addition of $r\in C^\infty$ cannot
remove any of these defects. This proves the last assertion as well.
\end{proof}

\begin{theorem}[Visibility criterion; \sourceT{39}]\label{thm:visibility}
Let
\[
 a_n=\int_M\Phi(x)^n\,d\mu(x),\qquad
 F(t)=\sum_{n\ge0}a_nt^n.
\]
For every nonzero same-side definite value $v$, the germ of $F$ has a continued
branch that fails to extend holomorphically through $t=1/v$. Consequently any
recurrence for $a_n$ over $\C(n)$ has shift order at least the number of distinct
nonzero same-side definite values of $\Phi$.
\end{theorem}
\begin{proof}
Since $M$ is compact, $\Phi$ is bounded and $F$ converges near zero. Define
\begin{equation}\label{eq:cauchy}
 G(z)=\int_M\frac{d\mu(x)}{z-\Phi(x)}
     =\int_\R\frac{d\nu(u)}{z-u}.
\end{equation}
The function $G$ is holomorphic in the upper half-plane. For large $|z|$,
\[
 G(z)=z^{-1}F(z^{-1}),\qquad F(t)=t^{-1}G(t^{-1}).
\]
Thus the upper-half-plane branch is connected to the germ at infinity obtained
from the original power series. It is not an unrelated solution of a
Picard--Fuchs equation.

We show that this branch of $G$ cannot extend holomorphically through $v$.
For $y>0$,
\begin{equation}\label{eq:poisson}
 -\frac1\pi\operatorname{Im}G(x+iy)
   =\int_\R\frac{1}{\pi}\frac{y}{(x-u)^2+y^2}\,d\nu(u).
\end{equation}
The right side is the Poisson approximation to the finite measure $\nu$.
It converges to $\nu$ in distributions as $y\downarrow0$.
For completeness, pair it with a compactly supported smooth test function
$\psi$: the inner integral is its Poisson average and converges uniformly to
$\psi(u)$, so dominated convergence against the finite measure proves the claim.

If the indicated branch of $G$ were holomorphic on a full disk centered at $v$,
its imaginary part on the real diameter would be real analytic. The convergence
on smaller diameters would then also be smooth, and \eqref{eq:poisson} would
force $\nu$ there to have the real-analytic density
$-\pi^{-1}\operatorname{Im}G(x)$. This contradicts Lemma~\ref{lem:density}.
The argument is valid even when $\rho$ is unbounded, because it uses
distributional boundary values instead of an unjustified pointwise formula at
$v$ itself.

Inversion carries a neighborhood of $v\ne0$ biholomorphically to a neighborhood
of $1/v$. Multiplication by $t^{-1}$ is harmless there. Hence the corresponding
branch of $F$ is genuinely singular at $1/v$. Distinct values give distinct
points, and Lemma~\ref{lem:budget} proves the recurrence lower bound.
\end{proof}

\begin{remark}[Why positivity alone would not suffice]\label{rem:cancellation}
A general critical point can be a saddle, and contributions on opposite sides
of the same critical value can have different regularity behavior. The proof
uses the stronger assertion that \emph{all} critical points on the fiber are
definite and have the same side. We will verify this by solving every nonzero
critical-point equation of the Franel torus. Checking one attractive critical
point while ignoring the others would not prove visibility.
\end{remark}

\begin{remark}[Non-global extrema]
The critical point need not maximize or minimize $\Phi$ on all of $M$.
The regular points lying over the same value contribute only a smooth density.
The continuation path from infinity runs in the upper half-plane; it does not
have to pass through intervening real critical values. These two observations
are what make interior critical values usable in the singularity count.
\end{remark}

\section{The Franel torus}\label{sec:torus}
Fix $s\ge2$ and put $d=s-1$. With $X=x_1\cdots x_d$, define the Laurent
polynomial
\begin{equation}\label{eq:Ps}
 P_s(x_1,\ldots,x_d)=\prod_{i=1}^{d}(1+x_i)(1+X^{-1}).
\end{equation}
All constant terms below are in the variables $x_1,\ldots,x_d$.

\begin{lemma}[Exact moment representation]\label{lem:CT}
For every $n\ge0$, $\CT P_s^n=A_s(n)$. On the unit torus $x_i=e^{i\theta_i}$,
$P_s$ is the real-valued function
\begin{equation}\label{eq:Phi}
 \Phi_s(\theta)=2^s\prod_{i=1}^{d}\cos(\theta_i/2)
                   \cos((\theta_1+\cdots+\theta_d)/2).
\end{equation}
Consequently, with normalized Haar measure on $\mathbb T^d$,
\[
 A_s(n)=\frac1{(2\pi)^d}\int_{\mathbb T^d}\Phi_s(\theta)^n\,d\theta,
 \qquad
 F_s(t)=\frac1{(2\pi)^d}\int_{\mathbb T^d}
              \frac{d\theta}{1-t\Phi_s(\theta)}.
\]
The generating-function identity holds for $|t|<2^{-s}$.
\end{lemma}
\begin{proof}
Expand each factor of $P_s^n$. If the term chosen from $(1+X^{-1})^n$ has
index $k$, its exponent in every variable is $-k$. A constant term requires
index $k$ from each $(1+x_i)^n$. The coefficient is then $\binom nk^s$, and
summing over $k$ proves the first assertion. Constant terms equal torus averages.
For the reality assertion use
$1+e^{i\theta}=2e^{i\theta/2}\cos(\theta/2)$. The phases from the first $d$
factors cancel the phase from the last one, giving \eqref{eq:Phi}.
Although individual half-angle cosines change sign under a $2\pi$ shift,
the product is periodic, as is also clear from \eqref{eq:Ps}.
Finally $|\Phi_s|\le2^s$, so the geometric series can be integrated uniformly
for the stated values of $t$.
\end{proof}

\begin{proposition}[All nonzero critical points]\label{prop:crit}
The critical points of $\Phi_s$ with nonzero value are exactly
\[
 (x_1,\ldots,x_d)=(\omega,\ldots,\omega),\qquad
 \omega^s=1,\quad\omega\ne-1.
\]
Their critical value is $(1+\omega)^s$.
\end{proposition}
\begin{proof}
At a point where $P_s\ne0$, no factor in \eqref{eq:Ps} vanishes. Logarithmic
derivatives, understood locally or simply as the rational derivative $P_s'/P_s$,
give
\begin{equation}\label{eq:logcrit}
 x_i\frac{\partial\log P_s}{\partial x_i}
       =\frac{x_i}{1+x_i}-\frac1{1+X}.
\end{equation}
On the real torus, $\partial/\partial\theta_i=i x_i\partial/\partial x_i$.
Since $P_s$ is real there, its real critical equations are precisely the
vanishing of \eqref{eq:logcrit}. Cross-multiplication gives $x_iX=1$ for each
$i$. Therefore all $x_i$ equal some $\omega$, and
$\omega X=\omega^{d+1}=\omega^s=1$.
Conversely each such point satisfies the equations. The nonzero requirement
excludes $\omega=-1$. Since $X^{-1}=\omega$ at the point, the value is
$(1+\omega)^{d+1}=(1+\omega)^s$.
\end{proof}

\begin{proposition}[Definiteness and critical values]\label{prop:hessian}
Let
\[
 m_s=\left\lfloor\frac{s+1}{2}\right\rfloor,\qquad
 v_j=(-1)^j\bigl(2\cos(\pi j/s)\bigr)^s,
 \quad 0\le j<s/2.
\]
These are exactly the distinct nonzero critical values of $\Phi_s$. Each is a
same-side definite value. At $x_i=e^{2\pi i j/s}$, the real Hessian is
\begin{equation}\label{eq:hess}
 \operatorname{Hess}_{\theta}\Phi_s
 =-(-1)^j\bigl(2\cos(\pi j/s)\bigr)^{s-2}(I_d+J_d),
\end{equation}
where $J_d$ is the all-ones matrix. There are $m_s$ such distinct values.
\end{proposition}
\begin{proof}
At a critical point, direct differentiation of \eqref{eq:logcrit} in
logarithmic coordinates gives
\[
 \operatorname{Hess}_{\log x}\log P_s
 =\frac{\omega}{(1+\omega)^2}(I_d+J_d).
\]
Indeed, the diagonal derivative of the first term in \eqref{eq:logcrit} is
$x_i/(1+x_i)^2$, while differentiating the second term with respect to any
$\log x_j$ gives $X/(1+X)^2$. At the critical point the latter equals
$\omega/(1+\omega)^2$ as well.
Passing to angular variables changes the sign. The gradient is zero, so
multiplication by $P_s=(1+\omega)^s$ gives the Hessian of $\Phi_s$, with no
additional gradient-square term. For $\omega=e^{2\pi i j/s}$,
\[
 \frac{\omega}{(1+\omega)^2}
    =\frac1{4\cos^2(\pi j/s)},\qquad
 (1+\omega)^s=(-1)^j\bigl(2\cos(\pi j/s)\bigr)^s.
\]
This proves \eqref{eq:hess}. The eigenvalues of $I_d+J_d$ are $1$ with
multiplicity $d-1$ and $d+1=s$ with multiplicity one. The Hessian is therefore
nondegenerate and definite. Positive $v_j$ are local maxima and negative
$v_j$ are local minima.

Complex conjugation pairs the indices $j$ and $s-j$, without changing the
value or Hessian side. The remaining representatives are exactly
$0\le j<s/2$. Their absolute values are strictly decreasing, since cosine is
strictly decreasing and positive on $[0,\pi/2)$. Consequently two different
representatives cannot give the same value. For odd $s$ there are $(s+1)/2$
representatives; for even $s$ there are $s/2$. Proposition~\ref{prop:crit}
excludes any additional critical point on a nonzero fiber, including a saddle.
All conditions of the visibility theorem have now been verified.
\end{proof}

\section{Proof of the minimum order}\label{sec:orderproof}
\begin{proof}[Proof of Theorem~\ref{thm:mainorder}]
For $s\ge2$, Lemma~\ref{lem:CT} supplies the real moment representation.
Proposition~\ref{prop:hessian} provides exactly $m_s$ distinct nonzero same-side
definite values. Theorem~\ref{thm:visibility} shows that the generating-function
germ has genuine continued singularities at all the points $1/v_j$.
Lemma~\ref{lem:budget} therefore gives $r\ge m_s$ for every recurrence of
shift order $r$, including those with coefficients in $\C(n)$ and finitely
many exceptional initial indices. Stoll's theorem \cite{Stoll} supplies a
recurrence over $\Q(n)$ of order at most $m_s$, proving equality.

For $s=1$, the binomial theorem gives $A_1(n)=2^n$ and the order-one relation
$A_1(n+1)-2A_1(n)=0$. An order-zero relation $p(n)A_1(n)=0$ with $p\ne0$
is impossible for all sufficiently large $n$. This proves the formula also
for $s=1$.
\end{proof}

\begin{example}[Small parameters]\label{ex:small}
For $s=2$ the only nonzero critical value is $4$, consistent with
$A_2(n)=\binom{2n}{n}$ and an order-one recurrence. For $s=3$ the values are
$8,-1$, so two distinct singularities force order at least two. For $s=4$ the
values are $16,-4$. For $s=5$ there are three values, and the two values apart
from $32$ are the roots of
\[
 x^2+11x-1=0.
\]
Their reciprocals are the roots of $t^2-11t-1=0$. This accounts for the
nonzero factors $(1-32t)(1+11t-t^2)$ in the highest-degree coefficient
polynomial of a corresponding differential equation, up to a nonzero scalar
and the convention for writing the shifts. These examples check the signs
and inversion in the general proof; they are not its evidence for all $s$.
\end{example}

\begin{audit}
The crucial additional content beyond an abstract singularity count is
Lemma~\ref{lem:density} and the exhaustive fiber calculation in
Propositions~\ref{prop:crit}--\ref{prop:hessian}. No claim that every critical
value of every real moment map is singular is used. No irreducibility of a
Picard--Fuchs operator is assumed. Nor does minimal order imply that every
Ap\'ery auxiliary solution satisfies a particular telescoping recurrence.
Those are different statements. External validation of this draft's general
minimality argument and its priority remains separate from the proof supplied
here.
\end{audit}

\section{Discrete Abel identities}\label{sec:abel}
We now turn to the independent arithmetic argument. A second-order recurrence
\begin{equation}\label{eq:general2}
 a_n f_{n+1}=b_n f_n+c_n f_{n-1},\qquad n\ge1,
\end{equation}
with $a_n\ne0$ is uniquely solved forward from two initial values. For two
solutions $U,f$, define
\[
 W_n(U,f)=U_nf_{n+1}-U_{n+1}f_n.
\]
The following computation is the discrete Abel identity; it is not a new
formula. Its use in Ap\'ery settings and the Franel specialization can be found
in \cite{CS,Tonien}.

\begin{lemma}[Discrete Abel identity]\label{lem:abel}
Two solutions of \eqref{eq:general2} satisfy
\[
 W_n=-\frac{c_n}{a_n}W_{n-1},\qquad
 W_n=(-1)^n\left(\prod_{j=1}^{n}\frac{c_j}{a_j}\right)W_0.
\]
\end{lemma}
\begin{proof}
Use \eqref{eq:general2} once for each solution:
\begin{align*}
 a_nW_n
 &=U_n(b_nf_n+c_nf_{n-1})-(b_nU_n+c_nU_{n-1})f_n\\
 &=c_n(U_nf_{n-1}-U_{n-1}f_n)=-c_nW_{n-1}.
\end{align*}
Iteration yields the product formula.
\end{proof}

\begin{theorem}[Prime-denominator rigidity]\label{thm:generalrigid}
Suppose a second-order recurrence has unique forward solutions and an integral
solution $U$ with $(U_0,U_1)\ne(0,0)$. Assume that for every solution $f$ its
Casoratian satisfies
\begin{equation}\label{eq:signature}
 W_n(U,f)=\frac{\kappa^n}{(n+1)^m}W_0(U,f),\qquad n\ge0,
\end{equation}
where $m\ge1$ and $\kappa\in\Z\setminus\{0\}$. Then every integral solution
is a rational scalar multiple of $U$. If $\gcd(U_0,U_1)=1$, it is an integer
scalar multiple.
\end{theorem}
\begin{proof}
For integral $f$, every $W_n$ is an integer. For a prime $p\nmid\kappa$, put
$n=p-1$ in \eqref{eq:signature}. The resulting equation
\[
 p^mW_{p-1}=\kappa^{p-1}W_0
\]
forces $p^m\mid W_0$. There are infinitely many such primes. The fixed integer
$W_0$ must therefore vanish. The two initial vectors are proportional over
$\Q$; uniqueness propagates that proportionality to the entire sequences.
If $(U_0,U_1)$ is primitive, choose integers $a,b$ with $aU_0+bU_1=1$.
For $f=qU$, the identity $q=af_0+bf_1$ then shows $q\in\Z$.
\end{proof}

\begin{remark}
Primitivity matters if the exact answer is written as $\Z U$. Without it,
the integral points on the rational line $\Q U$ can include nonintegral
multiples of the chosen $U$. The Franel normalization $S_0=1$ makes this
issue automatic. Also, two solutions with $W_0=0$ are identified through
initial-value uniqueness; dividing by terms $U_n$ that might vanish is
unnecessary in the general argument.
\end{remark}

\section{The cubic Franel lattice}\label{sec:lattice}
The binomial sum $S=A_3$ is integral and has $S_0=1,S_1=2$. It satisfies the
classical recurrence \eqref{eq:franel}, recorded in the original Franel work
and in \cite{Tonien,SZ}. For any other solution $f$, Lemma~\ref{lem:abel} gives
\begin{equation}\label{eq:franelW}
 W_n(S,f)=-\frac{8n^2}{(n+1)^2}W_{n-1}(S,f)
         =\frac{(-8)^n}{(n+1)^2}(f_1-2f_0).
\end{equation}
The product telescopes because
$\prod_{j=1}^n j^2/(j+1)^2=(n+1)^{-2}$.

\begin{proof}[Proof of Theorem~\ref{thm:mainlattice}]
Suppose first that $f$ is integral. For each odd prime $p$, equation
\eqref{eq:franelW} at $n=p-1$ implies
\[
 p^2\mid f_1-2f_0.
\]
Thus $f_1=2f_0$. The recurrence determines every later term, and the initial
vector is $f_0(1,2)$. Hence $f=f_0S$, with $f_0\in\Z$.
Conversely every integer multiple of $S$ is an integral solution.

For the second assertion, if $Df$ is integral, it is also a solution. The
first assertion gives $Df=aS$ for $a\in\Z$, so $f=(a/D)S$. Conversely a fixed
denominator for a rational scalar clears every term of $qS$.
\end{proof}

The recurrence is imposed for $n\ge1$, exactly as in the question. Were it
imposed at $n=0$ as well, it would immediately force $f_1=2f_0$, making the
question trivial. We do not add that condition. Instead, integrality over all
indices forces precisely that missing initial relation.

\begin{corollary}[No other integral ray]
Let $\mathcal V\subset\Q^{\N_0}$ be the solution space of
\eqref{eq:franel}. Then
\[
 \dim_\Q\mathcal V=2,\qquad
 \mathcal V\cap\Z^{\N_0}=\Z S,
\]
and the bounded-denominator subspace is $\Q S$.
\end{corollary}
\begin{proof}
The recurrence has nonzero forward coefficient $(n+1)^2$, so the initial-value
map $\mathcal V\to\Q^2$, $f\mapsto(f_0,f_1)$, is an isomorphism. The lattice
statements are Theorem~\ref{thm:mainlattice}.
\end{proof}

\section{Prime-adic strengthening and denominator certificates}\label{sec:padic}
The use of infinitely many primes can be replaced by arbitrarily large powers
of one fixed prime. This gives information that a mere common-denominator
classification does not display.

\begin{proposition}[Every odd prime detects a second direction]\label{prop:local}
Let $f$ be a rational solution of \eqref{eq:franel}, and put $c=f_1-2f_0$.
If $c\ne0$, then for every odd prime $p$ and every $a\ge1$,
\begin{equation}\label{eq:valuation}
 \min\{v_p(f_{p^a-1}),v_p(f_{p^a})\}\le v_p(c)-2a,
\end{equation}
where $v_p(0)=+\infty$. In particular, the denominators of $f_n$ have
unbounded powers of every odd prime.
\end{proposition}
\begin{proof}
At $n=p^a-1$, equation \eqref{eq:franelW} gives
\[
 v_p(W_{p^a-1})=v_p(c)-2a,
\]
since $p\nmid8$. Both $S_{p^a-1}$ and $S_{p^a}$ are integers. The
nonarchimedean triangle inequality therefore gives
\begin{align*}
 v_p(W_{p^a-1})
 &=v_p(S_{p^a-1}f_{p^a}-S_{p^a}f_{p^a-1})\\
 &\ge\min\{v_p(f_{p^a}),v_p(f_{p^a-1})\}.
\end{align*}
Rearranging proves \eqref{eq:valuation}. Its right side tends to $-\infty$.
For a rational number written in lowest terms, a negative valuation $-b$
means its denominator is divisible by $p^b$, proving the final assertion.
\end{proof}

\begin{corollary}[One-prime boundedness criterion]\label{cor:oneprime}
For a rational solution $f$, the following are equivalent:
\begin{enumerate}[label=(\roman*)]
\item $f\in\Q S$;
\item the sequence $(f_n)$ is bounded in $\Q_p$ for some odd prime $p$;
\item it is bounded in $\Q_p$ for every odd prime $p$.
\end{enumerate}
Also, if $\mathcal P$ is a finite set of primes, the solutions with all terms
in $\Z[\mathcal P^{-1}]$ are exactly
$\Z[\mathcal P^{-1}]S$.
\end{corollary}
\begin{proof}
For $f=qS$, integrality of $S$ gives $v_p(f_n)\ge v_p(q)$ uniformly.
Conversely Proposition~\ref{prop:local} excludes boundedness at even one odd
prime when $c\ne0$. For the localization statement, choose an odd prime outside
$\mathcal P$. Every element of the localization is integral at that prime, so
the one-prime criterion gives $f=qS$. The equality $q=f_0$ puts $q$ in the same
localization. The converse is immediate.
\end{proof}

\begin{example}[A concrete second solution]
Let $B$ solve \eqref{eq:franel} with $B_0=0$, $B_1=1$. Exact recurrence gives
\[
 B_2=4,\quad B_3=\frac{208}{9},\quad B_4=\frac{1280}{9}.
\]
Here $c=1$. The denominator $3^2$ is already visible at $n=3$. In general,
for every odd prime $p$, at least one of $B_{p-1},B_p$ has denominator divisible
by $p^2$, and the pair at $p^a-1,p^a$ contains a denominator divisible by
$p^{2a}$. Thus early integral values such as $B_1,B_2$ give no evidence for
integrality of the full sequence.
\end{example}

\begin{remark}[A finite certificate from initial values]
For a nonzero integer $c=f_1-2f_0$, any odd prime $p$ with $p^2\nmid c$
certifies that the entries cannot all be integral through index $p$. Indeed
$W_{p-1}=(-8)^{p-1}c/p^2$ is then nonintegral while its defining determinant
would be integral if all those entries were integral. This is a direct,
checkable certificate; it does not require the complete factorization of a
large later denominator.
\end{remark}

\section{Scope and provenance of the two arguments}
The source audit identifiers serve only to connect this paper with its
preparatory notes. They are not external theorem numbers.
\begin{center}
\begin{tabular}{@{}p{0.12\textwidth}p{0.35\textwidth}p{0.43\textwidth}@{}}
\toprule
Source & Result retained here & Essential justification\\
\midrule
T39 & Definite-value visibility & Positive same-side Morse density and
Poisson boundary values; every critical point in the fiber is covered.\\
T40 & General minimal shift order & Exhaustive torus critical-point and
Hessian calculation, visibility, singularity budget, Stoll upper bound.\\
T58 & Integral solution lattice & Discrete Abel identity and divisibility
of the one initial determinant by arbitrary odd-prime powers.\\
\bottomrule
\end{tabular}
\end{center}

The singularity budget is a standard recurrence-to-differential-equation
argument. The Morse lemma, the positivity of pushforward measures and Poisson
boundary convergence are classical analysis. The proposed contribution in
the first part is their complete specialization to the Franel torus, including
all nonzero critical fibers. The discrete Abel formula in the second part is
also classical. Its elementary prime-denominator consequence is what produces
the full integral-lattice classification and the stronger prime-adic statement.

The two outputs measure different kinds of rigidity. The minimum order says
that a sequence cannot be described by a shorter rational-coefficient relation.
The lattice statement says that one fixed two-dimensional rational recurrence
space has only one integral direction. Neither assertion is a transcendence
theorem for a named constant. In particular, no conclusion about irrationality
of an arbitrary Ap\'ery limit follows merely from unbounded denominators of a
second solution: such a conclusion also requires quantitative approximation
estimates of the appropriate strength.

No exhaustive priority certificate is asserted. The original general lower
bound is explicitly separated from the telescoping version in \cite{SZ};
Tonien records the remaining integral-sequence question in \cite{Tonien}.
The proofs above can therefore be evaluated against clearly specified existing
questions. Independent review must assess the argument and its place in the
literature rather than rely on the historical labels attached to the source
notes.

\appendix
\section{A compact proof audit}
This appendix identifies the exact points at which a tempting shorter argument
could fail, and the corresponding justification supplied in the paper.
\begin{enumerate}[label=\arabic*.]
\item \textbf{Reality of the moment map.} The complex phases in
\eqref{eq:Ps} cancel on the full torus. A generic Laurent polynomial would not
produce the real positive-measure Cauchy transform needed in
Theorem~\ref{thm:visibility}.
\item \textbf{Completeness of the critical list.} The derivation divides by
$P_s$ only on a nonzero fiber. It proves that every point on such a fiber is
either regular or one of the listed diagonal roots of unity. Degenerate
critical points at value zero do not enter the finite nonzero singularity
budget.
\item \textbf{No hidden saddle contribution.} Equation~\eqref{eq:hess} and
the exhaustive list imply definiteness of all critical points on the fiber.
Conjugate points contribute the same side and a positive leading density.
\item \textbf{Smoothness of the remaining density.} Compactness supplies a
whole interval on which the cutoff complement is a submersion. The proof
requires only a $C^\infty$ remainder, not an analytic partition of unity.
\item \textbf{A singularity of the correct branch.} The upper half-plane
connects the Cauchy transform at infinity to every real critical value from
above. Distributional Poisson convergence rules out a holomorphic extension
there. A local singularity of an unrelated period would be insufficient.
\item \textbf{Width rather than differential order.} The top differential
coefficient is $t^D H(t)$ with $\deg H\le r$, even when $D$ is much larger
than $r$. The number of nonzero singular points bounds $r$.
\item \textbf{Initial indices of the lattice question.} The recurrence starts
at $n=1$. The relation $f_1=2f_0$ is deduced from integrality and is not imposed
as an additional initial constraint.
\item \textbf{Fixed versus varying denominators.} Rationality of each term
alone is automatic and gives a two-dimensional space. A common denominator is
a genuinely stronger requirement. Proposition~\ref{prop:local} quantifies how
it fails away from the Franel line.
\end{enumerate}

\begin{thebibliography}{9}
\bibitem{Stoll}
M. Stoll, \emph{Bounds for the length of recurrence relations for convolutions
of P-recursive sequences}, European Journal of Combinatorics \textbf{18}
(1997), no.~6, 707--712.
\href{https://doi.org/10.1006/eujc.1996.0123}{doi:10.1006/eujc.1996.0123}.

\bibitem{SZ}
A. Straub and W. Zudilin, \emph{Sums of powers of binomials, their Ap\'ery
limits, and Franel's suspicions}, International Mathematics Research Notices
\textbf{2023} (2023), no.~11, 9861--9879.
\href{https://arxiv.org/abs/2112.09576}{arXiv:2112.09576}.

\bibitem{Tonien}
J. Tonien, \emph{Franel Numbers and a Continued Fraction Conjecture Discovered
by the Ramanujan Machine}, The Mathematical Intelligencer (2026), published
online 24 February 2026.
\href{https://doi.org/10.1007/s00283-025-10497-9}{doi:10.1007/s00283-025-10497-9}.

\bibitem{CS}
M. Chamberland and A. Straub, \emph{Ap\'ery limits: experiments and proofs},
American Mathematical Monthly \textbf{128} (2021), no.~9, 811--824.
\href{https://arxiv.org/abs/2011.03400}{arXiv:2011.03400}.
\end{thebibliography}
\end{document}
